%========================================================================================
% ENTORNO DE LA TESINA (ConTeXt LMTX)
%
% Sigue la guía de estilos (Guias/context_style_guide.md, copiada del repositorio de
% código): A4, márgenes de 2cm, cuerpo de 11pt, la paleta de colores, secciones en
% darkblue y tablas con encabezado gris. Agrega lo propio de una tesina: capítulos,
% páginas preliminares, índices, siglas, glosario y bibliografía. Lo carga tesina.tex
% con \environment entorno.
%========================================================================================

\startenvironment entorno

%----------------------------------------------------------------------------------------
% PLANTILLA BASE DE LA GUÍA DE ESTILOS
\setupcolors[state=start]
\setupcolor[hex]
\definecolor[darkblue][h=003366]
\definecolor[darkgray][h=e6e6e6]
\definecolor[redline][h=d9534f]
\definecolor[blueline][h=5bc0de]
\definecolor[greenline][h=5cb85c]

\setuphead[section][style=\bfa, color=darkblue, before=\blank, after=\blank]
\setuphead[subsection][style=\bf, color=black, before=\blank, after=\blank]
\setuphead[subsubsection][style=\bf, color=black, before=\blank, after=\blank]
\setuppapersize[A4][A4]
\setupbodyfont[11pt]
\setuplayout[width=middle, height=middle, topspace=2cm, bottomspace=2cm]

%----------------------------------------------------------------------------------------
% IDIOMA Y PÁRRAFOS
\mainlanguage[es]
\language[es]

% Separación entre párrafos equivalente a un \blank, sin sangría.
\setupwhitespace[big]

%----------------------------------------------------------------------------------------
% ENLACES Y MARCADORES DEL PDF (enlaces sin color)
\setupinteraction
  [state=start,
   color=,
   contrastcolor=,
   style=,
   title={Optimización picking y ruteo de pedidos en un centro de distribución},
   author={Ximena Ebertz, Lucía Soria}]

\placebookmarks[preliminar,title,chapter,section,subsection][force=yes]
\setupinteractionscreen[option=bookmark]

%----------------------------------------------------------------------------------------
% NUMERACIÓN DE PÁGINAS
% Romanos en las páginas preliminares y arábigos desde el Capítulo 1 (tesina.tex
% reinicia el número ahí).
\defineconversionset[frontpart:pagenumber][][romannumerals]
\defineconversionset[bodypart:pagenumber][][numbers]
\defineconversionset[appendix:pagenumber][][numbers]
\defineconversionset[backpart:pagenumber][][numbers]
\setupuserpagenumber[numberconversionset=pagenumber]

%----------------------------------------------------------------------------------------
% ETIQUETAS
\setuplabeltext
  [es]
  [chapter=Capítulo~,
   appendix=Anexo~,
   figure=Figura~,
   table=Cuadro~]

%----------------------------------------------------------------------------------------
% CAPÍTULOS Y PÁGINAS PRELIMINARES
% La guía no define capítulos: cada uno empieza en una página nueva, con
% "Capítulo N" (o "Anexo A") en una línea y el título debajo, en negro.
\protected\def\TesinaTituloCapitulo#1#2%
  {\vbox
     {\hsize\textwidth
      \setupalign[flushleft,nothyphenated]%
      \doifsomething{#1}{#1\par\vskip.5\lineheight}%
      #2}}

\setuphead
  [chapter]
  [page=yes,
   command=\TesinaTituloCapitulo,
   numberstyle=\bfb,
   textstyle=\bfc,
   after={\blank[2*big]}]

\setuphead
  [title]
  [page=yes,
   command=\TesinaTituloCapitulo,
   textstyle=\bfc,
   after={\blank[2*big]},
   incrementnumber=list]

% Páginas preliminares con título (resumen, listas, índices, dedicatoria).
\definehead[preliminar][title]

%----------------------------------------------------------------------------------------
% ÍNDICES
\setupcombinedlist
  [content]
  [list={preliminar,title,chapter,section}]

\setuplist
  [preliminar]
  [style=normal,
   pagestyle=normal,
   before=]

\setuplist
  [title,chapter]
  [style=bold,
   pagestyle=bold,
   before={\blank[medium]}]

%----------------------------------------------------------------------------------------
% TABLAS
% Estilo de la guía: \bTABLE[split=yes, align=middle, frame=on, width=fit] con la
% primera fila de encabezado en gris y negrita:
%     \setupTABLE[row][1][background=color, backgroundcolor=darkgray, style=bold]

%----------------------------------------------------------------------------------------
% SIGLAS Y GLOSARIO
%
% \ac{tag}: la primera vez escribe "forma larga (sigla)" y después solo la sigla.
% Las listas muestran todo lo definido acá (van antes del texto, así que no pueden
% filtrar por uso).
\definesynonyms[sigla][siglas][\siglalarga][\siglacorta]
\setupsynonyms[sigla][headstyle=bold, location=serried, width=fit, distance=1em, criterium=all]

\definesynonyms[termino][terminos][\terminodefinicion][\terminonombre]
\setupsynonyms[termino][headstyle=bold, location=serried, width=fit, distance=1em, criterium=all]

\protected\def\ac#1%
  {\ifcsname tesina:ac:#1\endcsname
     \siglacorta{#1}%
   \else
     \global\expandafter\let\csname tesina:ac:#1\endcsname\relax
     \siglalarga{#1} (\siglacorta{#1})%
   \fi}

% Abreviaciones
\sigla[mip]{MIP}{Programación Lineal Entera Mixta (del inglés {\em Mixed-Integer Programming})}
\sigla[hc] {HC} {Ascenso de Colina (del inglés {\em Hill Climbing})}
\sigla[aco]{ACO}{Optimización por Colonia de Hormigas (del inglés {\em Ant Colony Optimization})}

% Glosario
\termino[pickeo]            {Picking}            {Proceso de recolección de los ítems que componen un pedido dentro de un almacén, también llamado preparación de pedidos}
\termino[picker]            {Picker}             {Operario encargado de recolectar los ítems de uno o más lotes dentro del almacén}
\termino[lote]              {Lote}               {Conjunto de pedidos agrupados que se asignan y recolectan de forma conjunta}
\termino[pedido]            {Pedido}             {Solicitud de un cliente, compuesta por uno o más ítems a recolectar}
\termino[deposito]          {Depósito}           {Nodo de partida y llegada del grafo que representa el almacén (también llamado {\em depot})}
\termino[duedate]           {Due date}           {Fecha de entrega comprometida de un pedido, respecto de la cual se penaliza el adelanto o la demora en su recolección}
\termino[batching]          {Batching}           {Agrupamiento de pedidos en lotes}
\termino[cadenadefrio]      {Cadena de frío}     {Restricción de tiempo máximo que un pedido puede permanecer fuera de temperatura controlada durante su recolección}
\termino[instancia]         {Instancia}          {Conjunto particular de datos de entrada (almacén, pedidos, lotes, parámetros) sobre el que se ejecuta un modelo}
\termino[escenariofuncional]{Escenario funcional}{Escenario de prueba de tamaño y complejidad moderados, pensado para validar la correctitud de un modelo}
\termino[escenariodeestres] {Escenario de estrés}{Escenario de prueba de mayor escala, pensado para evaluar el comportamiento de un modelo en condiciones exigentes}

%----------------------------------------------------------------------------------------
% BIBLIOGRAFÍA
% Estilo numérico (aps), ordenado alfabéticamente por autor, con el DOI visible y como
% enlace.
\usebtxdataset[default][Referencias/tesina.bib]
\usebtxdefinitions[aps]
\setupbtx[dataset=default]
\setupbtxrendering[aps][numbering=yes, sorttype=short]

\starttexdefinition mutable protected btx:aps:doi-url #1
    #1
    \btxdoif {doi} {
        \btxcomma
        \goto {DOI:~\hyphenatedurl{\btxflush{doi}}} [url(https://doi.org/\btxflush{doi})]
    }
\stoptexdefinition

%----------------------------------------------------------------------------------------
% FLOTANTES PARA PSEUDOCÓDIGOS
\definefloat[algoritmo][algoritmos]
\setuplabeltext[es][algoritmo=Algoritmo~]

\stopenvironment
